\chapter{Gamification Hub}
\label{ch:hubteoria}
La vecchia istanza polyglot è spenta (HTTP 530): l'integrazione obbligatoria
è la \textbf{nuova Hub} di UniAQ.

\section{Componenti}
Console: \url{https://gamification-webapp.createlab-univaq.it};
API: \url{https://gamification-api.createlab-univaq.it/api/v1};
Swagger: \url{https://gamification-api.createlab-univaq.it/swagger-ui/index.html};
guida: \url{https://gamification-api.createlab-univaq.it/docs/api-guide.en.md};
video: \url{https://youtu.be/EJ8l2kcWiFs}.

\section{Modello concettuale}
Autenticazione \texttt{POST /auth \{username, password, origin:''GAME''\}}
$\rightarrow$ Bearer JWT valido 24\,h. Runtime
\texttt{POST /executions \{gameId, playerId, actionId, data:\{\}\}} con
\texttt{data} obbligatorio. Concetti: \textit{actions} (eventi di gioco),
\textit{point concepts} con periodi in millisecondi, \textit{badge} come
collezioni, un solo \textit{level} per point concept, \textit{regole Drools}
con \texttt{/rules/validate} e simulazioni, \textit{challenge/leaderboard/
classifiche}.

\section{Migrazione}
I vecchi \texttt{gameId} non sono validi: il gioco va ricreato da console,
le regole riscritte in Drools e validate via Swagger prima del go-live.
Questo sposta su Hub proprio ciò che Toda segnala come gap (leaderboard,
economia XP/energia) senza riscrivere il client (si veda \texttt{specifica.pdf},
Cap.~6).
